% Custom preamble for Charles Coverdale's R Journal-style papers.
% Drop into the paper/rj/ directory and reference via the Rmd YAML:
%   output:
%     rticles::rjournal_article:
%       includes:
%         in_header: header.tex

% Tables
\usepackage{longtable}
\usepackage{tabularx}
\usepackage{booktabs}

% Float placement: use floating [tbp] rather than hard [H] for better
% page packing. Tune LaTeX's defaults to allow bigger fractions of a
% page to be given to a float before it drifts to a dedicated page.
\usepackage{float}
\usepackage{placeins}  % for \FloatBarrier
\renewcommand{\topfraction}{0.92}
\renewcommand{\bottomfraction}{0.8}
\renewcommand{\floatpagefraction}{0.7}
\renewcommand{\textfraction}{0.07}
\setcounter{topnumber}{3}
\setcounter{bottomnumber}{2}
\setcounter{totalnumber}{4}

% Leave page bottoms ragged rather than stretching vertical space
% between figures or sections to justify to the bottom margin.
% Essential for papers with many figures: without this, LaTeX opens
% 2-3 inch gaps between stacked floats to reach the bottom line.
\raggedbottom

% Suppress the auto-generated \address{...} trailer. The R Journal
% template emits an author block from the YAML metadata at the end of
% the paper; Charles's preference is title page byline only.
\renewcommand{\address}[1]{}

% Numbered sections and subsections. pandoc maps markdown ## to
% \subsection, which RJournal.sty leaves unnumbered by default.
% Override the titleformat so subsections display a single arabic
% counter ("1 Introduction", "2 Background", ...).
\setcounter{secnumdepth}{3}
\usepackage{titlesec}
\renewcommand\thesubsection{\arabic{subsection}}
\renewcommand\thesubsubsection{\arabic{subsection}.\arabic{subsubsection}}
\titleformat{\subsection}{\normalfont\large\bfseries}{\thesubsection}{1em}{}
\titleformat{\subsubsection}{\normalfont\normalsize\bfseries}{\thesubsubsection}{1em}{}

% Blue DOI / URL / citation hyperlinks (RJournal.sty leaves colorlinks off).
\hypersetup{colorlinks=true,
  linkcolor=blue!60!black,
  urlcolor=blue!60!black,
  citecolor=blue!60!black}

% Make DOIs in the bibliography clickable. natbib / plainnat prints
% "doi: 10.xxxx" as plain text; redefine \doi so the string is wrapped
% in \href and gets the blue colorlinks treatment.
\renewcommand{\doi}[1]{\href{https://doi.org/#1}{doi:~#1}}
